
\documentclass{article}
\usepackage{colortbl}
\usepackage{makecell}
\usepackage{multirow}
\usepackage{supertabular}

\begin{document}

\newcounter{utterance}

\centering \large Interaction Transcript for game `wordle\_withclue', experiment `medium\_frequency\_words\_clue\_no\_critic', episode 3 with Qwen3.5{-}9B.
\vspace{24pt}

{ \footnotesize  \setcounter{utterance}{1}
\setlength{\tabcolsep}{0pt}
\begin{supertabular}{c@{$\;$}|p{.15\linewidth}@{}p{.15\linewidth}p{.15\linewidth}p{.15\linewidth}p{.15\linewidth}p{.15\linewidth}}
    \# & \multicolumn{2}{c}{Player} && \multicolumn{2}{c}{Game Master} \\
    \hline

    \theutterance \stepcounter{utterance}  
    & & & \multicolumn{4}{p{0.6\linewidth}}{
        \cellcolor[rgb]{0.9,0.9,0.9}{
            \makecell[{{p{\linewidth}}}]{
                \texttt{\tiny{[P1$\langle$GM]}}
                \texttt{You are a language wizard who likes to guess words by using the given rules.} \\
\\ 
\texttt{Welcome to Wordle! You have six attempts to guess the target word, a valid English word of five lowercase letters (a{-}z). Please use the tags "explanation:" and "guess:" to provide a concise explanation for each guess.} \\
\\ 
\texttt{To help you make an informed guess, you will receive a clue for the word, such as} \\
\texttt{clue: snowy white.} \\
\\ 
\texttt{Here is an example guess based on the clue:} \\
\texttt{explanation: In the fairy tail Snow White, the girl is killed because she eats a poisoned apple. And the word apple has 5 letters.} \\
\texttt{guess: apple} \\
\\ 
\texttt{After each guess, your answer will be validated, and you will receive feedback indicating which letters are correct (green), which letters are correct but in the wrong position (yellow), and which letters are incorrect (red). This feedback can be useful in determining which letters to include or exclude in your next guess.} \\
\\ 
\texttt{For example, the feedback for "apple" might be:} \\
\texttt{guess\_feedback: a<yellow> p<yellow> p<green> l<yellow> e<red>} \\
\\ 
\texttt{The explanation should contain details about how both the guess\_feedback and clue work together in tandem to formulate a new guess.} \\
\\ 
\texttt{Do not generate any other text.} \\
\\ 
\texttt{Let's begin with your first guess.} \\
\\ 
\texttt{clue: sweetheart} \\
            }
        }
    }
     \\ \\

    \theutterance \stepcounter{utterance}  
    & \multicolumn{4}{p{0.6\linewidth}}{
        \cellcolor[rgb]{0.561,0.737,0.561}{
            \makecell[{{p{\linewidth}}}]{
                \texttt{\tiny{[P1$\rangle$GM]}}
                \texttt{explanation: The clue "sweetheart" suggests a word related to affection or a partner. "SWEET" is a direct synonym for the first part of the clue and is a common five{-}letter word. It also contains common vowels and consonants that might appear in the target word.} \\
\texttt{guess: sweet} \\
            }
        }
    }
    & & \\ \\

    \theutterance \stepcounter{utterance}  
    & & & \multicolumn{2}{p{0.3\linewidth}}{
        \cellcolor[rgb]{0.9,0.9,0.9}{
            \makecell[{{p{\linewidth}}}]{
                \texttt{\tiny{[GM$|$GM]}}
                \texttt{attempts = 1} \\
\texttt{target\_word = flame} \\
\texttt{target\_word\_clue = sweetheart} \\
\texttt{guess = sweet} \\
\texttt{guess\_feedback = s<red> w<red> e<yellow> e<red> t<red>} \\
            }
        }
    }
    & & \\ \\

    \theutterance \stepcounter{utterance}  
    & & & \multicolumn{4}{p{0.6\linewidth}}{
        \cellcolor[rgb]{0.9,0.9,0.9}{
            \makecell[{{p{\linewidth}}}]{
                \texttt{\tiny{[P1$\langle$GM]}}
                \texttt{guess\_feedback: s<red> w<red> e<yellow> e<red> t<red>} \\
\\ 
\texttt{Provide your response only in this format.} \\
\texttt{explanation: details} \\
\texttt{guess: word} \\
            }
        }
    }
     \\ \\

    \theutterance \stepcounter{utterance}  
    & \multicolumn{4}{p{0.6\linewidth}}{
        \cellcolor[rgb]{0.561,0.737,0.561}{
            \makecell[{{p{\linewidth}}}]{
                \texttt{\tiny{[P1$\rangle$GM]}}
                \texttt{explanation: The feedback indicates that 'e' is present but in the wrong position, while 's', 'w', and 't' are incorrect. Since the clue "sweetheart" implies a word related to love or affection, and we know 'e' is in the word but not in the second or fourth position, I will try "heart" related words or synonyms. However, "heart" is only 5 letters but doesn't fit the 'e' position well if it's not the first. Let's try "lover" or "dear". Wait, "dear" is 4 letters. How about "sweet" was wrong on s,w,t. Maybe the word is "heart"? No, h,e,a,r,t. e is second. Feedback said e is yellow (wrong position). So e is not second. What about "sweet" clue? Maybe "sugar"? s,u,g,a,r. s is red. No. Maybe "honey"? h,o,n,e,y. e is fourth. Feedback said e is not fourth (sweet had e in 3rd and 4th, both yellow/red? Wait, sweet: s(1) red, w(2) red, e(3) yellow, e(4) red, t(5) red. So e is in the word, but not at pos 3 or 4. So e is at 1, 2, or 5. Let's try "heart" again? h,e,a,r,t. e is 2. If e is yellow in sweet (pos 3), it means e is in the word but not at 3. In heart, e is at 2. That's allowed. But s,w,t are not in heart. So heart is possible. But clue is sweetheart. Maybe "lover"? l,o,v,e,r. e is 4. But feedback said e in sweet (pos 3 and 4) was yellow/red. Wait, sweet has two e's. Pos 3 and 4. Feedback: e<yellow> e<red>. So one e is correct letter wrong pos, one is wrong letter? Or both are correct letter but positions? Usually Wordle feedback per letter. So first e (pos 3) is yellow (correct letter, wrong pos). Second e (pos 4) is red (wrong letter). So e is in the word, but not at pos 3 or 4. And s,w,t are not in the word. So we need a 5{-}letter word with e not at 3 or 4, no s,w,t. Clue sweetheart {-}> maybe "honey"? h,o,n,e,y. e is 4. But e cannot be at 4. So no. "sugar" has s. No. "candy"? c,a,n,d,y. No e. "choco"? No} \\
            }
        }
    }
    & & \\ \\

    \theutterance \stepcounter{utterance}  
    & & & \multicolumn{2}{p{0.3\linewidth}}{
        \cellcolor[rgb]{0.9,0.9,0.9}{
            \makecell[{{p{\linewidth}}}]{
                \texttt{\tiny{[GM$|$GM]}}
                \texttt{Error: The response should contain the 'guess:' keyword exactly once.} \\
            }
        }
    }
    & & \\ \\

    \theutterance \stepcounter{utterance}  
    & & & \multicolumn{2}{p{0.3\linewidth}}{
        \cellcolor[rgb]{0.9,0.9,0.9}{
            \makecell[{{p{\linewidth}}}]{
                \texttt{\tiny{[GM$|$GM]}}
                \texttt{game\_result = ABORT} \\
            }
        }
    }
    & & \\ \\

\end{supertabular}
}

\end{document}
